\section{Contracting the base while preserving adjoint rank}
\label{g:cf:contraction}

A relative section is useful only if it survives the next change of base.
The theorem below makes that requirement explicit. The section records
existence on a complete generic fiber; the one-dimensional adjoint space
controls the ratios that arise when a larger fiber is formed by contraction.
The proof constructs two different metrics on the base. The first is
bounded and controls the relative section. The second comes from the
pushforward volume and controls cotangent tensors. Their common numerical
zero directions determine the contraction.

\begin{theorem}[Rank-preserving contraction]\label{g:cf:contraction-theorem}
Let $X$ be a smooth connected projective complex variety with
$H^0(X,\Omega_X^q)=0$ for $q>0$, and suppose that $L=-K_X$ has a smooth
semipositive Hermitian metric. Let an algebraic torus $T$ act on $X$,
with its natural action on $L$. Let $f:X\dashrightarrow Z$ be a dominant
$T$-invariant rational map to a smooth connected projective variety,
where $\mathbb C(Z)$ is relatively algebraically closed in $\mathbb C(X)$.
On a smooth projective resolution $\pi:Y\to X$, $f_Y:Y\to Z$, put
$L_Y=\pi^*L$, $E_Y=K_{Y/X}$, and
$\eta=\operatorname{Spec}\mathbb C(Z)$. Suppose that:
\begin{enumerate}
\item $H^0(Y_\eta,pL_Y|_{Y_\eta})^T\ne0$ for some integer $p>0$;
\item there is an integer $a>0$ such that
\[
\dim_{\mathbb C(Z)}H^0(Y_\eta,(E_Y+mL_Y)|_{Y_\eta})^T=1
\quad\text{for every positive multiple }m\text{ of }a.
\]
\end{enumerate}
Then either $H^0(X,rL)^T\ne0$ for some $r>0$, or there is a dominant
rational map $h:Z\dashrightarrow W$ to a smooth connected projective
variety with $\dim W<\dim Z$ such that $h\circ f$ satisfies the same two
conditions, possibly with different positive integers. The field
$\mathbb C(W)$ is relatively algebraically closed in $\mathbb C(X)$.
\end{theorem}

We prove the theorem in four steps. Vertical orders normalize the given
section and identify the base primes at which poles disappear on $X$.
Fiber integration then supplies the bounded metric and a nef signed
log-anticanonical class. A supporting movable class produces effective
divisors in unbounded twists. Finally, interpolation of their horizontal
orders gives the new section, and the same field of section ratios restores
adjoint rank one.

\subsection{The two generic-fiber conditions}
\label{g:cf:generic-conditions}

Average the weights of the given metric over a maximal compact subgroup of
$T$.  The resulting smooth semipositive metric $h$ is invariant under that
compact subgroup.  Via the identity section of $K_X+L$, it defines a smooth
positive volume form $d\mu_X$.

Fix the rational map and resolution in Theorem~\ref{g:cf:contraction-theorem},
writing $f:Y\to Z$ for the resolved map. We call its two hypotheses the
\emph{section condition} and the \emph{rank condition}. The action on the
base and its bundles is trivial; $E_Y=K_{Y/X}$ has its natural invariant
Jacobian section.

Both conditions are unchanged by resolving further or by replacing the
base birationally.  Indeed, if $\nu:Y'\to Y$ is a further smooth
modification, then
\[
 E_{Y'}=\nu^*E_Y+K_{Y'/Y},
 \qquad \nu_*\mathcal O_{Y'}(K_{Y'/Y})=\mathcal O_Y.
\]
The comparison of the adjoint spaces uses multiplication by the relative
Jacobian section.  The same comparison on a common resolution defines the
invariant rational sections even on a model on which the torus action is
not regular.  The degree-zero rank is a consequence of both hypotheses: multiplication by a
suitable power of the section in the first condition gives an injection
into a positive-degree invariant space of dimension one (choose a common
multiple of its degree and the integer in the rank condition), while the
Jacobian section gives a nonzero element at degree zero.

If $\dim Z=0$, the section condition gives a global section on $Y$, and
then on $X$.  We may therefore assume $d=\dim Z>0$.  Every smooth base
encountered below has no positive-degree holomorphic forms: pullback
under a dominant map is injective, and holomorphic forms are invariant
under smooth projective birational modification.  Consequently
\begin{equation}\label{g:cf:euler}
 \chi(Z,\mathcal O_Z)=1.
\end{equation}
Our task is to obtain a global invariant section or construct another
map satisfying both conditions with a smaller base.

\subsection{Equidimensional preparation and the exceptional pole divisor}
\label{g:cf:preparation}

Apply weak toroidalization to a resolution of the rational map,
prescribing the complement of an isomorphism open over $X$ as part of the
source boundary \cite[Theorem~1.1]{ADK2013}.  Then apply the equidimensional
modification by projective subdivisions
\cite[Proposition~4.4, arXiv:alg-geom/9707012v1]{AK2000}.  After renaming the base, we have
\[
 \pi_0:Y_0\longrightarrow X,
 \qquad f_0:Y_0\longrightarrow Z,
\]
where $Y_0$ is normal and toroidal, $Z$ is smooth projective, $\pi_0$ is
projective birational, and $f_0$ is projective, surjective, equidimensional,
and toroidal.  The reduced source boundary $\Sigma$ contains every
$\pi_0$-exceptional divisor.  The base boundary has simple normal
crossings.  This step is birational: it does not assert that the source
is smooth or that the fibers are reduced.  We may subsequently use
smooth toroidal resolutions by subdivisions \cite{KKMSD}; these remain
toroidal over $Z$ but need not remain equidimensional.

Let $s$ be a rational section supplied by the section condition.  On
$Y_0$, its horizontal divisor is effective.  Write
\begin{equation}\label{g:cf:normalization}
 \frac1p\operatorname{div}(s)=H+f_0^*V,
\end{equation}
where the coefficient of $V$ at a prime $P\subset Z$ is the minimum of
\[
 \frac{\operatorname{coeff}_D(\operatorname{div}(s)/p)}
      {\operatorname{mult}_D(f_0^*P)}
 \quad(D\text{ a component of }f_0^*P).
\]
Equidimensionality ensures that every vertical prime maps onto a prime
of $Z$.  Thus $V$ is a rational divisor and $H$ is effective and
$\mathbb Q$-Cartier.  Choose $r>0$ sufficiently divisible.  There is an
invariant rational section
\[
 \sigma\in H^0\bigl(Y_0,r\pi_0^*L-f_0^*(rV)\bigr),
 \qquad \operatorname{div}(\sigma)=rH.
\]
Its pullback to every dominating smooth model is regular. We take $r$
also divisible by the integer in the rank condition.  For every
prime of the base there is at least one component above it with
coefficient zero in $H$.

Let $G$ be the reduced sum of the primes $P\subset Z$ for which every
component of $f_0^*P$ having coefficient zero in $H$ is exceptional over
$X$.  This is a finite sum, since each such prime is the image of an
exceptional divisor of $\pi_0$.  Every nonexceptional component above
$G$ has strictly positive coefficient in $H$.  Hence, for every
sufficiently small rational $\delta>0$,
\begin{equation}\label{g:cf:exceptional-negative}
 (H-\delta f_0^*G)^-
 \text{ is exceptional over }X.
\end{equation}
If $V+\delta G$ is rationally effective for such a $\delta$, multiply its
section by a suitable power of $\sigma$ and divide by the corresponding
canonical section of $\delta G$.  Equation
\eqref{g:cf:exceptional-negative} shows that the resulting invariant
rational section of a positive multiple of $L$ has no pole along any
prime of $X$.  Smoothness of $X$ then extends it globally.  Thus it
suffices either to prove this effectivity or to contract the base.

\subsection{The bounded metric and the pushforward density}
\label{g:cf:metrics}

Write
\begin{equation}\label{g:cf:rho}
 \rho(z)=\frac{f_*(\pi^*d\mu_X)}{d\lambda(z)}
\end{equation}
on the smooth locus, using any smooth resolution of the prepared
toroidal model. The next estimate uses only the smoothness and
positivity of $d\mu_X$; its curvature and the adjoint ranks enter
afterwards, when the density is made into a metric.

Write the base boundary as $\sum_iP_i$.  For a component $D$
above $P_i$, let
\[
 b_{iD}=\operatorname{mult}_D(f_0^*P_i),\qquad
 e_D=\operatorname{coeff}_D(K_{Y_0}-\pi_0^*K_X),
\]
where canonical divisors are chosen compatibly using a rational top
form on $X$.  At the regular generic point of $D$, $e_D$ is its
Jacobian order; it is nonnegative, and zero when $D$ is
nonexceptional.  Define
\begin{equation}\label{g:cf:thresholds}
 c_i=\min_{D\mapsto P_i}\frac{1+e_D}{b_{iD}}>0,
 \qquad B=\sum_i(1-c_i)P_i=B^+-B^-.
\end{equation}
If $c_i>1$, every component above $P_i$ is exceptional, so $P_i$ is
a component of $G$.  Thus $B^-$ is supported on the primes all of whose components
are exceptional over $X$, and in particular on $G$.  The coefficients of $B^+$ are strictly less than one.


\begin{lemma}[Toroidal density bounds]\label{g:cf:density-bounds}
For the prepared toroidal diagram and smooth positive volume above,
near an arbitrary base point,
with $y_i=0$ the local boundary branches, for every $\eta>0$,
\begin{equation}\label{g:cf:upper-density}
 \rho(y)\leq C_\eta
       \prod_i|y_i|^{-2(1-c_i+\eta)}
\end{equation}
on the general smooth locus.  Near a general point of $P_i$,
\begin{equation}\label{g:cf:lower-density}
 \rho(y)\geq C|y_i|^{-2(1-c_i)},\qquad C>0.
\end{equation}
At a general point of a prime outside the base boundary, $\rho$ also
has a positive constant lower bound.

\end{lemma}

\begin{proof}
The divisor
\begin{equation}\label{g:cf:effective-log-jacobian}
 K_{Y_0}+\Sigma-\pi_0^*K_X-\sum_i c_i f_0^*P_i
\end{equation}
is effective and $\mathbb Q$-Cartier.  Indeed, no prime contributes to
the pullbacks of two distinct $P_i$, by equidimensionality.  At a
component above $P_i$, its coefficient is $1+e_D-c_i b_{iD}\geq0$.
At the other boundary components it is $1+e_D$, and away from the
boundary it is $e_D$.  The divisor $K_{Y_0}+\Sigma$ is Cartier and
pulls back log crepantly under toroidal subdivisions: the local top
logarithmic form of a toric model trivializes it and pulls back to the
top logarithmic form on the subdivision.  Consequently, on a smooth
toroidal resolution, if $x_j=0$ are the local source boundary branches,
then
\begin{equation}\label{g:cf:orders}
 a_j:=1+e_j>0,\qquad
 a_j\geq\sum_i c_i b_{ij},
\end{equation}
where $e_j$ is the Jacobian order over $X$ and $b_{ij}$ is the order
of $y_i$ along $x_j=0$.  These inequalities include every new ray on every toroidal subdivision,
in particular the exceptional rays over intersections of boundary
components. They follow from pullback of an effective log-Cartier divisor,
so this is more than an estimate at the original boundary primes.

For use in analytic integration we make explicit why formal toroidal
charts can be replaced by analytic monomial charts.  In a formal chart
write
\begin{equation}\label{g:cf:formal-units}
 f^*y_i=\prod_jx_j^{A_{ij}},\qquad
 x_j=X_jU_j,\quad y_i=Y_iV_i.
\end{equation}
In \eqref{g:cf:formal-units}, the indices include all source and
target toric coordinates, including unit coordinates; thus
$A$ is a $d\times\dim X$ integer matrix. Here $X_j$ and $Y_i$ are
analytic boundary equations, or the constant
$1$ for unit coordinates; $U_j,V_i$ are formal units.  The coordinates
minus their residues are regular parameters.  Dominance implies
that $A$ has full row rank, since the logarithmic monomial map has
generic rank equal to the target dimension.  Factoring the analytic
boundary equations gives analytic units $R_i$ with
\[
 f^*Y_i=\Bigl(\prod_jX_j^{A_{ij}}\Bigr)R_i.
\]
Choose analytic units $V_i'$ having the same first jets as $V_i$.
We seek analytic units $U_j'$ satisfying
\begin{equation}\label{g:cf:unit-equations}
 \prod_j(U_j')^{A_{ij}}=R_i f^*V_i'.
\end{equation}
Use the constants of the formal solution, divide by them, and take
logarithms near $1$.  The resulting equations are linear equations
with matrix $A$ and holomorphic right hand sides.  A linear right
inverse gives a holomorphic solution.  Its first jet differs from
the prescribed formal first jet by a vector in $\ker A$; a
kernel-valued affine correction removes that difference without
changing the equations.  Exponentiating gives the desired $U_j'$.
Thus $X_jU_j'$ and $Y_iV_i'$ have the original first jets and remain
coordinate systems, while the monomial identities hold exactly
analytically.  Negative exponents can occur only on unit
coordinates and cause no difficulty.

We can now estimate by coarea in logarithmic coordinates.  Cut
arguments into bounded sectors and work on relatively compact
pieces of these charts, away from the boundary itself.  Relative
to Lebesgue measure in the logarithmic variables, the pullback
of $d\mu_X$ is bounded above by a constant times
\[
 \prod_j|x_j|^{2a_j}
\]
over the vanishing coordinates; all unit moduli are bounded above
and away from zero.  The factor $|x_j|^2$ in this formula is the
Jacobian of the logarithmic change of variables.  The analogous
change on the base introduces $\prod_i|y_i|^2$, up to bounded
positive factors from its unit coordinates.  The logarithmic
map is the fixed full-rank complex linear map $A$, up to boundedly
many argument translates.  Its pushforward density is therefore
an integral over affine real planes with fixed coarea factors.

In the following products, $i$ ranges only over vanishing target
coordinates and $j$ only over vanishing source coordinates; $b_{ij}$
is the corresponding boundary-order block of $A$.
To estimate $\rho\prod_i|y_i|^{2(1-c_i+\eta)}$, divide the
logarithmic pushforward density by
$\prod_i|y_i|^{2(c_i-\eta)}$.  On an affine fiber, this factor is
constant and can be brought into the integrand.  The resulting
exponent on a vanishing source coordinate is
\begin{equation}\label{g:cf:residual-exponents}
 2\Bigl(a_j-\sum_i(c_i-\eta)b_{ij}\Bigr)>0.
\end{equation}
If some $b_{ij}>0$, strict positivity follows from
\eqref{g:cf:orders} and $\eta>0$; otherwise it follows from $a_j>0$.
The integrand thus decays exponentially as the real logarithms
of vanishing coordinates tend to $-\infty$.  Their upper bounds,
and the bounds on all remaining logarithmic coordinates, allow
us to dominate it by $C\exp(-\varepsilon\|u\|)$ on the whole
chart piece in real logarithmic variables.  Integrals of this
bound over affine planes of a fixed direction are uniformly
bounded independently of the translation.  For example, write
such a plane as $v+\ker A$ with $v$ orthogonal to $\ker A$ and
use $\|v+w\|\geq\|w\|$.  Properness supplies finitely many
chart pieces over a sufficiently small base neighborhood.  The
coordinate equations of each boundary prime differ only by
bounded units on overlaps.  Summing the bounds proves
\eqref{g:cf:upper-density} almost everywhere, and hence on the
general locus where the density is continuous.

For \eqref{g:cf:lower-density}, take a general point of a component
which realizes the minimum in \eqref{g:cf:thresholds}.  At that point take coordinates
\begin{equation}\label{g:cf:power-chart}
 y_1=x_1^{b_{iD}},\qquad y_j=x_j\quad(2\leq j\leq d).
\end{equation}  The pulled-back volume is
comparable there to $|x_1|^{2e_D}$ times ordinary coordinate
volume.  The change $y_1=x_1^{b_{iD}}$, followed by integration
over a fixed patch in the remaining fiber variables, contributes
\[
 |y_1|^{2((1+e_D)/b_{iD}-1)}
   =|y_1|^{-2(1-c_i)}.
\]
This is a lower bound for the full pushforward density.

At a general point of an excluded prime outside the base boundary,
the density also has a positive constant lower bound.  In fact
there is a nonexceptional component over that prime.  An
exceptional component would lie in the source boundary, whereas
a boundary component meeting the inverse image of the base
toroidal open cannot be vertical there: all target coordinates
are units, so the full-rank monomial matrix uses only source
unit coordinates and the map restricted to the component is
dominant.  At a nonexceptional component the Jacobian order is
zero. If $b$ is its pullback multiplicity over the base prime,
the same power-coordinate calculation gives a lower bound,
since its exponent is $2(1/b-1)\leq0$.
\end{proof}

\begin{lemma}\label{g:cf:metric-conclusions}
In the preceding setup, $V$ has a semipositive metric with locally bounded
weights. For the effective rational divisors $B^+$ and $B^-$ defined
in \eqref{g:cf:thresholds}, there is a nef rational divisor $M$ such that
\begin{equation}\label{g:cf:log-relation}
 K_Z+B^++M\sim_{\mathbb Q}B^-.
\end{equation}
The support of $B^+$ has simple normal crossings and all its coefficients
are less than one.  The divisor $B^-$ is supported on those primes of $Z$ all of whose
components on $Y_0$ are exceptional over $X$; in particular it is
supported on $G$.
\end{lemma}

\begin{proof}
We first match the data with the smooth invariant version of
Lemmas~\ref{g:moments} and~\ref{g:transfer}. Their original variety is
$X$, their allowed-pole divisor is zero, and both input bundles are
$A=L=-K_X$ with the compact-invariant smooth metric $h$. The finite
measure is $d\mu_X$. Take exponent $r$ and the rational section
$\widetilde s$ of $rL$ whose pullback has divisor $rH+f_0^*(rV)$.
It is the power of the original generic section used to construct
$\sigma$. It is invariant and has no horizontal poles. On an equivariant
resolution the invariant adjoint ranks equal one for every multiple of
$r$, including zero by the argument above. All hypotheses of
Lemma~\ref{g:moments} are therefore satisfied.

The zeroth density is $\rho$ from \eqref{g:cf:rho}; hence
$-\log\rho$ is psh on the smooth base-change locus. The normalized
base divisor in Lemma~\ref{g:transfer} is $rV$, by the full-component
minimum in \eqref{g:cf:normalization}. Its bounded metric, divided by
$r$, gives the asserted metric on $V$. More concretely, for a local frame
$v$ of $rV$, it is determined by
\[
 |v|^2=\left(\max_{f_0^{-1}(z)}
      |\sigma f_0^*v|_{(\pi_0^*h)^r}^{2/r}\right)^r.
\]
The essential supremum in the transfer lemma is this maximum: on a
smooth connected general fiber the smooth measure has full support.
The invariance step is the orthogonal quotient in the smooth
Berndtsson theorem, as proved in Lemma~\ref{g:moments}; no singular
invariant-subbundle positivity is invoked.

We now use Lemma~\ref{g:cf:density-bounds} to construct the signed
log-anticanonical metric. Put $M=-K_Z-B$.  On the general open its local metric weight is
\begin{equation}\label{g:cf:corrected-weight}
 \varphi_M=-\log\rho-\sum_i(1-c_i)\log|y_i|^2.
\end{equation}
This has the transition law of $M$: a frame of $-K_Z$ is
multiplied by the boundary factor $\prod_i y_i^{1-c_i}$,
after clearing denominators.  It is plurisubharmonic away
from the boundary.  The lower density bounds give the upper
bound needed for extension across a dense open of every
excluded divisor.  Hartogs extension treats the remaining
codimension-two set.  The upper density estimate implies
\begin{equation}\label{g:cf:arbitrarily-small-poles}
 \varphi_M\geq\eta\sum_i\log|y_i|^2-O_\eta(1)
 \qquad\text{for every }\eta>0.
\end{equation}
In particular this metric has zero Lelong numbers: the lower estimate
holds with arbitrarily small logarithmic coefficients. For every fixed
divisible $k>0$, choose $\eta<1/k$.
Normal-crossing integrability shows that
$\mathcal I(k\varphi_M)=\mathcal O_Z$.  This assertion is
about every such power, not just the first multiplier ideal.

Choose a very ample divisor $A$ on $Z$.  Nadel vanishing
\cite{Nadel}, applied after adding the smooth positive metric
on $(d+1-i)A$, gives
\[
 H^i\bigl(Z,K_Z+kM+(d+1-i)A\bigr)=0
 \qquad(1\leq i\leq d).
\]
Castelnuovo--Mumford regularity makes
$K_Z+kM+(d+1)A$ globally generated.  Its class divided by $k$
is nef, and the limit as $k$ tends to infinity proves that
$M$ is nef.  The identical argument for the bounded metric
on $V$, whose every power has trivial multiplier ideal,
also shows that $V$ is nef.  Equation
\eqref{g:cf:log-relation} follows from the definition of $M$.
\end{proof}

\subsection{Cotangent tensors on a supporting movable class}
\label{g:cf:slopes}

If $V+G$ is big, then for $0<\delta\leq1$,
\[
 V+\delta G=\delta(V+G)+(1-\delta)V
\]
is big, because $V$ is pseudoeffective.  Choosing a sufficiently
small rational $\delta$ and using
\eqref{g:cf:exceptional-negative} finishes the proof.  We may
therefore assume that $V+G$ is not big.  Duality between the
pseudoeffective divisor cone and the closed movable curve cone
\cite{BDPP2013} gives a nonzero movable curve class $\beta$ with
\[
 (V+G)\cdot\beta=0.
\]
Since $V$ and $G$ are pseudoeffective and $B^-$ is supported
on $G$, we obtain
\begin{equation}\label{g:cf:null-degrees}
 V\cdot\beta=G\cdot\beta=B^-\cdot\beta=0.
\end{equation}

\begin{lemma}\label{g:cf:cotangent-slope}
Every line subsheaf $N\subseteq(\Omega_Z^1)^{\otimes s}$, for
$s\geq0$, satisfies $c_1(N)\cdot\beta\leq0$.
\end{lemma}

\begin{proof}
For $s=0$, the inclusion into $\mathcal O_Z$ gives the
assertion directly.  Let $s>0$ and let $t>0$ be rational.
Since $M$ is nef, $M+tA$ is ample.  Choose a sufficiently
divisible very ample multiple and a general member; divide
that member by the multiple and add it to $B^+$ to obtain a
rational simple normal crossings boundary $\Delta_t$, all
of whose coefficients lie in $[0,1)$.  Then
\[
 K_Z+\Delta_t\sim_{\mathbb Q}B^-+tA
\]
is pseudoeffective.

Apply the orbifold cotangent theorem
\cite[Theorem~1.3 and Section~5]{CampanaPaun2019} to this smooth
projective log-canonical pair.  On a finite adapted cover
$\tau:\widetilde Z\to Z$, the adapted orbifold cotangent
bundle $\mathcal E_t$ has determinant class
$\tau^*(K_Z+\Delta_t)$.  Every torsion-free quotient of each
$\mathcal E_t^{\otimes s}$ has nonnegative degree against
$\tau^*\beta$.  This holds for classes on the boundary of
the movable cone as well.

There is an inclusion $\tau^*\Omega_Z^1\subseteq\mathcal E_t$.
In codimension one, at a branch $z=w^b$ corresponding to
boundary coefficient $\gamma$, the orbifold generator is
$w^{-b\gamma}dz$; hence the ordinary generator $dz$ belongs
to the adapted bundle.  Additional branch divisors have
$\gamma=0$.  Let $\mathcal N$ be the saturation of the
pulled-back line $\tau^*N$ in $\mathcal E_t^{\otimes s}$.
Saturation adds an effective divisorial class, so
\[
 \deg(\tau)\,c_1(N)\cdot\beta
   \leq c_1(\mathcal N)\cdot\tau^*\beta.
\]
The quotient by $\mathcal N$ is torsion-free.  Its determinant
has nonnegative degree, and the tensor determinant identity
for a bundle of rank $d$ is
\[
 c_1(\mathcal E_t^{\otimes s})
      =s d^{s-1}c_1(\mathcal E_t).
\]
If the quotient has rank zero, it is zero and the same
inequality holds with determinant degree zero.  Dividing
by the cover degree and using \eqref{g:cf:null-degrees}
gives
\begin{equation}\label{g:cf:tensor-bound}
 c_1(N)\cdot\beta
   \leq s d^{s-1}(K_Z+\Delta_t)\cdot\beta
   =s d^{s-1}tA\cdot\beta.
\end{equation}
Letting $t$ decrease to zero proves the assertion.
\end{proof}

\begin{lemma}\label{g:cf:effective-sequence}
There are a fixed line bundle $N$ and effective divisors $D_j$
such that
\begin{equation}\label{g:cf:divisor-sequence}
 D_j\sim m_jV+N,\qquad
 m_j\longrightarrow+\infty,\qquad
 D_j\cdot\beta=0,
\end{equation}
where the positive integers $m_j$ may be chosen strictly increasing
and sufficiently divisible.
\end{lemma}

\begin{proof}
Choose $b>0$ with $bV$ Cartier. Lemma~\ref{g:cf:metric-conclusions}
gives $bV$ bounded semipositive weights, and the base has no
positive-degree forms by the argument preceding \eqref{g:cf:euler}.
Apply Lemma~\ref{g:determinant} to $bV$, prescribing any further
divisibility required for the construction. It gives a cotangent line
$N$ and effective divisors $D_j\sim N+n_j(bV)$ with $n_j\to\infty$.
Set $m_j=bn_j$. Lemma~\ref{g:cf:cotangent-slope} gives
$N\cdot\beta\leq0$, while $V\cdot\beta=0$. Effectivity gives the
opposite inequality for $D_j\cdot\beta$, so every degree is zero.
\end{proof}

\subsection{Maximal zero-degree systems and horizontal interpolation}
\label{g:cf:zero-systems}

Choose a linear system of divisor class $D$ with
$D\cdot\beta=0$ whose rational map has maximal possible
image dimension among all such systems.  The constant
map is permitted.  This maximal dimension is smaller
than $d$: a generically finite system makes its divisor
big, whereas a nonzero supporting functional cannot
vanish on a big class.  Resolve this rational map and
take its Stein factorization, followed by a smooth
projective resolution of the Stein base.  We obtain
\[
 h:\widehat Z\longrightarrow Z,\qquad
 \widehat g:\widehat Z\longrightarrow W,
 \qquad \dim W<d,
\]
where the general fibers of $\widehat g$ are connected.

\begin{lemma}\label{g:cf:maximal-field}
Every ratio of two nonzero sections in any one linear
system whose divisor class has degree zero against
$\beta$ belongs to $\mathbb C(W)\subseteq\mathbb C(Z)$.
\end{lemma}

\begin{proof}
Multiply the sections of that system by those defining
the chosen maximal system.  Their products form a
system of degree zero, and its map is the product of
the two maps followed by the Segre embedding.  By
maximality it has the same image dimension as the
chosen map.  The added section ratios are therefore
algebraic over the field of that image.  The Stein
factorization identifies $\mathbb C(W)$ with the
relative algebraic closure of the image field in
$\mathbb C(Z)$.  This proves the assertion.
\end{proof}

The interpolation below is a relative form of the divisor argument in
\cite[proof of Theorem~3.1, Step~3]{LP2018}. There, zero Iitaka
dimension makes linearly equivalent effective divisors equal. Here
the maximal-field property gives equality of their horizontal parts.
The subsequent rank-restoration argument still has to control their
vertical poles.

\begin{lemma}\label{g:cf:horizontal-slope}
The rational divisor $h^*V$ is rationally effective
on the complete generic fiber of $\widehat g$.
\end{lemma}

\begin{proof}
Use the sequence in \eqref{g:cf:divisor-sequence}.
For $j<k<l$, linear equivalence gives
\begin{equation}\label{g:cf:interpolation}
 (m_l-m_j)D_k
   \sim (m_l-m_k)D_j+(m_k-m_j)D_l.
\end{equation}
Both sides are effective divisors in the same
zero-degree system.  By Lemma~\ref{g:cf:maximal-field},
their ratio is a rational function from $W$.
Its pullback has no horizontal divisor over $W$.
Thus \eqref{g:cf:interpolation}, after pullback by
$h$, is an equality of actual horizontal divisors.

Fix the first two indices.  For every horizontal
prime $Q$ on $\widehat Z$, the coefficients of
$h^*D_j$ have the form $a_Q+m_jb_Q$.  The preceding
identity proves this simultaneously for every
index.  Since all $D_j$ are effective and
$m_j\to\infty$, every $b_Q$ is nonnegative.
These slopes have finite support: they are the
coefficients of the horizontal part of
\[
 \frac{h^*(D_2-D_1)}{m_2-m_1}
     \sim_{\mathbb Q}h^*V.
\]
This representative of $h^*V$ has effective
horizontal part.  Its vertical components disappear
on the generic fiber, yielding the stated linear
effectivity there.  The argument uses actual
divisors and rational functions, rather than only
numerical equivalence.
\end{proof}

\subsection{Preserving the section and rank conditions}
\label{g:cf:rank-preservation}

Compose $X\dashrightarrow Z$ with $Z\dashrightarrow W$.
The composition is invariant and has connected general
fibers: the inclusions of function fields remain
relatively algebraically closed.  Choose a smooth
projective resolution $\pi':Y'\to X$ dominating $Y_0$
and the map to $\widehat Z$, and write
\[
 f':Y'\longrightarrow\widehat Z,\qquad
 u=\widehat g\circ f'.
\]
Multiplication of the pullback of $\sigma$ by the
generic section of $h^*V$ from
Lemma~\ref{g:cf:horizontal-slope} gives an invariant
section of a positive multiple of $L_{Y'}$ on the
generic fiber of $u$.  Thus the section condition
is preserved.

To prove the rank condition, fix a sufficiently
divisible $m>0$ and take any two nonzero invariant
sections of $E_{Y'}+mL_{Y'}$ over the generic fiber
of $u$.  They can be extended simultaneously to
global sections of
\begin{equation}\label{g:cf:spreading-twist}
 E_{Y'}+mL_{Y'}+ku^*A_W
\end{equation}
for a sufficiently large integer $k$, where $A_W$
is very ample.  For completeness, every pole
divisor has image in a proper closed subset of
$W$.  A high multiple of a very ample bundle
has a section in any prescribed power of the
ideal of this finite union of images.  Pullback
then vanishes to at least the required orders
on all the finitely many pole divisors.  The
same section clears the poles of both rational
sections.  This also handles images of
codimension greater than one.  It is a base
section, so it preserves invariance and does
not change their ratio.

Divide the two sections of
\eqref{g:cf:spreading-twist} by the Jacobian
section and $\sigma^{m/r}$.  The old rank
condition says that the results come from
rational sections on $\widehat Z$ of
\[
 mh^*V+k\widehat g^*A_W.
\]
Push their divisors down to $Z$.  A pole outside
$G$ is impossible: for a prime $P\not\subseteq G$,
choose a component above $P$ on $Y_0$ with
coefficient zero in $H$ and nonexceptional over
$X$.  Its strict transform on $Y'$ has
coefficient zero in both the pulled-back $H$
and the Jacobian divisor $E_{Y'}$.  Since $h$
is an isomorphism at the generic point of $P$,
a pole downstairs would remain a pole on this
component, contradicting regularity of the
original global section.

After bounding the poles of both sections
along $G$ by a common integer $\ell$, their
ratio is a ratio of sections in the single
linear system
\begin{equation}\label{g:cf:rank-system}
 mV+k h_*\widehat g^*A_W+\ell G.
\end{equation}
We verify that its degree against $\beta$ is
zero.  Only the middle term needs proof.
If $W$ is a point it is zero.  Otherwise,
let $A_I$ be the pullback to $W$ of the
hyperplane class of the image of the
defining maximal system.  The map from
$W$ to that image is generically finite,
so $A_I$ is big.  On a sufficiently high
resolution of the system,
\[
 h^*D\sim\widehat g^*A_I+F,
 \qquad F\geq0.
\]
Some sufficiently large multiple of $A_I$
minus $A_W$ is rationally effective, by
bigness.  Pulling back and pushing forward
shows that a positive multiple of $D$
minus $h_*\widehat g^*A_W$ is rationally
effective.  The latter divisor is itself
effective after choosing an effective
representative of $A_W$.  Since
$D\cdot\beta=0$ and movable classes are
nonnegative on effective divisors, both
inequalities force
\[
 (h_*\widehat g^*A_W)\cdot\beta=0.
\]
Together with \eqref{g:cf:null-degrees},
this proves the asserted zero degree of
\eqref{g:cf:rank-system}.

Lemma~\ref{g:cf:maximal-field} now puts
the ratio of the two original sections
in $\mathbb C(W)$.  Thus the new
invariant generic adjoint space has
dimension at most one for every positive multiple of a common integer.  The new section
condition and the Jacobian supply a
nonzero element in every sufficiently
divisible degree, giving dimension
exactly one.  Both conditions have
therefore been preserved while the
base dimension has strictly decreased.

This proves Theorem~\ref{g:cf:contraction-theorem}. The construction is
conditional on the two generic-fiber hypotheses, and their preservation is
part of its conclusion. In particular the argument can be iterated without
introducing a strict-curvature assumption on the base.
